%%%%%%%%%%%%%%%%%%%%%%%%%%%%%%%%%%%%%%%%%%%%%%%%%%%%%%%%%%%%%%%%%%%%%%%%%%%%%%%%%%%%%%%%%%
% Analysis introduction 
%%%%%%%%%%%%%%%%%%%%%%%%%%%%%%%%%%%%%%%%%%%%%%%%%%%%%%%%%%%%%%%%%%%%%%%%%%%%%%%%%%%%%%%%%%

NEEDS TO BE REWRITTEN FOR NEW TOPOLOGY.


The signal is searched for in the phase space defined by the kinematic variables $m_{ES}$ and $\DeltaE$. 
This provides clean identification of the $B$ decay as these events occupy a 
limited section of this phase space. We will use data outside of the signal region to determine the background shape of the fit. \\

The signal shape parameters will be extracted from signal Monte Carlo events run through the \babar ~{} GEANT simulation. \\

\subsection{Datasets}

BaBar collected data from 1998 to 2008 and the data was broken into 6 run periods, each representing
roughly a year of datataking. In Table~\ref{tab:dataskims} we show the integrated luminosity, 
the number of events in the skim (discussed in the next section), the number of events before
the skim, and the estimated number of $B\bar{B}$ pairs, calculated from the cross section for
$e^+e^- \rightarrow B\bar{B}$.


